\section{Dynamic Modeling \& Small-Signal Analysis (Week 10)}

\subsection{Keywords \& Comprehensive Overview}
\noindent\textbf{Ctrl+F keywords:} dynamic modeling, small-signal analysis, state-space averaging, SSA, perturbation, linearization, quiescent operating point, DC steady state, AC dynamic equation, state vector, input vector, switch ON matrix, switch OFF matrix, averaged matrix, resolvent matrix, characteristic equation, characteristic polynomial, control-to-output transfer function, $G_{vd}(s)$, line-to-output transfer function, audio susceptibility, $G_{vg}(s)$, control-to-inductor current, $G_{id}(s)$, output impedance, $Z_{out}(s)$, input impedance, $Z_{in}(s)$, Buck converter, Boost converter, Buck-Boost converter, Right-Half-Plane zero, RHP zero, non-minimum phase, phase lag, undershoot, loop crossover frequency, bandwidth limitation, phase margin, Bode plot, second-order canonical form, natural frequency, corner frequency, quality factor, damping ratio, capacitor ESR zero.

\vspace{0.5em}
\noindent Switched-mode power converters are fundamentally non-linear, time-varying dynamic systems because circuit topologies change abruptly as semiconductor switches toggle between ON and OFF states. Steady-state volt-second and capacitor charge balance principles predict only the DC quiescent operating point ($X = [I_L, V_o]^T$). However, real-world power electronic systems are subject to continuous dynamic disturbances, including:
\begin{itemize}
  \item Step changes and ripples in the input line voltage ($v_{in}(t) = V_{in} + \hat{v}_{in}(t)$).
  \item Sudden load resistance and load current steps ($i_o(t) = I_o + \hat{\imath}_o(t)$).
  \item Start-up transients and reference tracking commands ($v_{ref}(t)$).
  \item Closed-loop feedback controller design to ensure stability, fast settling time, and adequate phase and gain margins.
\end{itemize}
To apply classical and modern frequency-domain control theory (Bode plots, Nyquist criterion, root locus, PID and lead-lag compensator design), the non-linear switching converter must be modeled as a continuous, linear time-invariant (LTI) small-signal transfer function. The \textbf{State-Space Averaging (SSA)} method solves this by:
\begin{enumerate}
  \item \textbf{Averaging:} Removing high-frequency switching ripples by time-weighting the piecewise state equations over the switching period $T_s$, valid whenever the switching frequency is much higher than the converter natural resonant frequency ($f_s \gg f_0$).
  \item \textbf{Perturbation \& Linearization:} Decomposing all variables into large-signal DC quiescent values plus infinitesimal AC perturbations ($x = X + \hat{x}$, $d = D + \hat{d}$), and neglecting second-order product terms ($\hat{d}\hat{x} \approx 0$).
\end{enumerate}

\subsubsection{Master Summary Table: Small-Signal State-Space Models of DC-DC Converters}
\begin{center}
\footnotesize
\renewcommand{\arraystretch}{1.25}
\begin{tabularx}{\textwidth}{@{} >{\raggedright\arraybackslash}p{2.0cm} >{\raggedright\arraybackslash}p{2.6cm} >{\raggedright\arraybackslash}p{4.8cm} >{\raggedright\arraybackslash}p{2.6cm} >{\raggedright\arraybackslash}X @{}}
\toprule
\textbf{Topology} & \textbf{Averaged ($A, B, F$)} & \textbf{Transfer Function $G_{vd}(s)$} & \textbf{Poles \& Zeros} & \textbf{Dynamic Behavior \& Stability} \\
\midrule
\textbf{Buck} \newline (Step-Down) &
$A = \begin{bmatrix} 0 & -\frac{1}{L} \\ \frac{1}{C} & -\frac{1}{RC} \end{bmatrix}$ \newline
$B = \begin{bmatrix} D/L \\ 0 \end{bmatrix}, \; F = \begin{bmatrix} V_{in}/L \\ 0 \end{bmatrix}$ &
$G_{vd}(s) = \dfrac{V_{in}}{L C s^2 + \dfrac{L}{R} s + 1}$ \newline
$= \dfrac{V_{in}}{1 + \dfrac{s}{Q\omega_0} + \dfrac{s^2}{\omega_0^2}}$ &
$\omega_0 = \dfrac{1}{\sqrt{LC}}$ \newline
$Q = R\sqrt{\dfrac{C}{L}}$ \newline
No RHP zero (minimum phase). &
\textbf{Minimum Phase.} Phase lag $\to -180^\circ$. Easily compensated with Type-II/III voltage-mode control. High bandwidth achievable ($f_c \approx f_s/5$). \\
\midrule
\textbf{Boost} \newline (Step-Up) \newline ($D' \equiv 1-D$) &
$A = \begin{bmatrix} 0 & -\frac{D'}{L} \\ \frac{D'}{C} & -\frac{1}{RC} \end{bmatrix}$ \newline
$B = \begin{bmatrix} 1/L \\ 0 \end{bmatrix}, \; F = \begin{bmatrix} V_o/L \\ -I_L/C \end{bmatrix}$ &
$G_{vd}(s) = \dfrac{V_o}{D'} \dfrac{1 - s \dfrac{L}{{D'}^2 R}}{1 + \dfrac{L}{{D'}^2 R} s + \dfrac{LC}{{D'}^2} s^2}$ &
$\omega_0 = \dfrac{D'}{\sqrt{LC}}$ \newline
$Q = D' R \sqrt{\dfrac{C}{L}}$ \newline
$\omega_{z,RHP} = \dfrac{{D'}^2 R}{L}$ &
\textbf{Non-Minimum Phase (RHP Zero).} Transient undershoot before rising. Adds $+20\,\text{dB/dec}$ gain but $-90^\circ$ phase lag! Severe bandwidth limit: $f_c < \frac{1}{3 \sim 5} f_{z,RHP}$. \\
\midrule
\textbf{Buck-Boost} \newline (Inverting) \newline ($D' \equiv 1-D$) &
$A = \begin{bmatrix} 0 & \frac{D'}{L} \\ -\frac{D'}{C} & -\frac{1}{RC} \end{bmatrix}$ \newline
$B = \begin{bmatrix} 1/L \\ 0 \end{bmatrix}, \; F = \begin{bmatrix} \frac{V_{in} - V_o}{L} \\ \frac{I_L}{C} \end{bmatrix}$ &
$G_{vd}(s) = \dfrac{V_o}{D D'} \dfrac{1 - s \dfrac{D L}{{D'}^2 R}}{1 + \dfrac{L}{{D'}^2 R} s + \dfrac{LC}{{D'}^2} s^2}$ &
$\omega_0 = \dfrac{D'}{\sqrt{LC}}$ \newline
$Q = D' R \sqrt{\dfrac{C}{L}}$ \newline
$\omega_{z,RHP} = \dfrac{{D'}^2 R}{D L}$ &
\textbf{Non-Minimum Phase (RHP Zero).} Inductor disconnected from output when switch is ON. Phase lag reaches $-270^\circ$. Bandwidth heavily limited under heavy load ($R$ small). \\
\bottomrule
\end{tabularx}
\end{center}

---

\subsection{State-Space Averaging (SSA) 6-Step Method}

\subsubsection{Theoretical Rationale: Why Steady-State Analysis is Insufficient}
In DC steady state, switching converters satisfy:
\begin{enumerate}
  \item \textbf{Inductor Volt-Second Balance:} $\int_0^{T_s} v_L(t)\,dt = 0 \implies \langle v_L \rangle_{T_s} = 0$.
  \item \textbf{Capacitor Charge Balance:} $\int_0^{T_s} i_C(t)\,dt = 0 \implies \langle i_C \rangle_{T_s} = 0$.
\end{enumerate}
These algebraic balances determine only the constant DC operating points ($V_o = D V_{in}$, $I_L = I_o$, etc.). However, when the duty cycle $d(t)$ is dynamically modulated by a feedback controller, or when the input line voltage $v_{in}(t)$ or load current fluctuates, the instantaneous net energy stored in $L$ and $C$ is non-zero over any given switching cycle:
\begin{equation}
L \frac{d\langle i_L \rangle}{dt} = \langle v_L(t) \rangle_{T_s} \ne 0, \qquad C \frac{d\langle v_C \rangle}{dt} = \langle i_C(t) \rangle_{T_s} \ne 0
\end{equation}
State-Space Averaging rigorously links the circuit's differential equations in each switching sub-interval to its dynamic transfer functions.

\subsubsection{The Formal 6-Step SSA Procedure}

\paragraph{Step 1: Define State Vector and Input Vector}
Choose independent energy storage variables as the state vector $x(t)$, and the external independent sources as the input vector $u(t)$:
\begin{equation}
x(t) = \begin{bmatrix} i_L(t) \\ v_C(t) \end{bmatrix}, \qquad u(t) = v_{in}(t) \quad (\text{or } [v_{in}(t),\, i_o(t)]^T)
\end{equation}
The desired system output (typically load voltage $v_o(t)$) is expressed via an output equation:
\begin{equation}
y(t) = v_o(t) = C_m x(t) + E_m u(t)
\end{equation}

\paragraph{Step 2: Sub-interval 1 (Switch ON, $0 < t < d T_s$)}
During the interval $d T_s$, the active switch is conducting and the diode is reverse biased. Formulate the circuit differential equations using Kirchhoff's Voltage Law (KVL) for each inductor and Kirchhoff's Current Law (KCL) for each capacitor node:
\begin{equation}
\dot{x}(t) = A_1 x(t) + B_1 u(t) \qquad (\text{or } A_{on} x + B_{on} u)
\end{equation}
where $A_1 \in \mathbb{R}^{2\times 2}$ is the state matrix and $B_1 \in \mathbb{R}^{2\times 1}$ is the input matrix during the ON state.

\paragraph{Step 3: Sub-interval 2 (Switch OFF, $d T_s < t < T_s$)}
During the complementary interval $(1-d)T_s$, the switch is open and the freewheeling diode conducts. Formulate the corresponding state equations:
\begin{equation}
\dot{x}(t) = A_2 x(t) + B_2 u(t) \qquad (\text{or } A_{off} x + B_{off} u)
\end{equation}
where $A_2 \in \mathbb{R}^{2\times 2}$ and $B_2 \in \mathbb{R}^{2\times 1}$ represent the OFF-state circuit dynamics.

\paragraph{Step 4: Time-Averaging over One Switching Period $T_s$}
Because the ON and OFF states persist for fractions $d(t)$ and $1-d(t)$ of each period respectively, compute the continuous moving average of the state derivative over $T_s$:
\begin{equation}
\dot{x}(t) = \left[ d(t) A_1 + (1 - d(t)) A_2 \right] x(t) + \left[ d(t) B_1 + (1 - d(t)) B_2 \right] u(t)
\end{equation}
\begin{itemize}
  \item \textbf{Averaging Approximation Validity:} The moving average accurately represents the system dynamics under the condition that the switching frequency is substantially higher than the system natural resonance frequency and modulation frequency ($f_s \gg f_0$). Under this condition, the high-frequency triangular switching ripple is effectively filtered out by the $LC$ low-pass filter, leaving only the low-frequency dynamic trajectory.
  \item \textbf{Non-linearity:} Notice that this averaged state equation is non-linear because the control input $d(t)$ multiplies the state vector $x(t)$ (product $d(t)x(t)$).
\end{itemize}

\paragraph{Step 5: Small-Signal Perturbation around DC Operating Point}
Assume small perturbations around a steady-state quiescent operating point:
\begin{equation}
x(t) = X + \hat{x}(t), \qquad u(t) = V_{in} + \hat{v}_{in}(t), \qquad d(t) = D + \hat{d}(t)
\end{equation}
where:
\begin{itemize}
  \item $X = \begin{bmatrix} I_L \\ V_C \end{bmatrix}$, $V_{in}$, and $D$ denote constant DC equilibrium values ($d/dt = 0$).
  \item $\hat{x}(t) = \begin{bmatrix} \hat{\imath}_L(t) \\ \hat{v}_C(t) \end{bmatrix}$, $\hat{v}_{in}(t)$, and $\hat{d}(t)$ represent small AC variations.
  \item \textbf{Small-Signal Condition:} The perturbations are infinitesimal relative to the DC values:
  \begin{equation}
  \|\hat{x}\| \ll \|X\|, \qquad |\hat{v}_{in}| \ll V_{in}, \qquad |\hat{d}| \ll D \text{ and } |\hat{d}| \ll (1 - D)
  \end{equation}
\end{itemize}

\paragraph{Dr. Liu's Numerical Justification for Neglecting Second-Order Products (Lecture Slide 31):}
Expanding the non-linear product of perturbed quantities yields:
\begin{equation}
(D + \hat{d})(I + \hat{\imath}) = D I + D \hat{\imath} + I \hat{d} + \hat{d} \hat{\imath}
\end{equation}
Consider realistic operating values:
\begin{itemize}
  \item Nominal DC values: $D = 0.5$, $I = 2.0\,\si{\ampere}$.
  \item Small-signal AC perturbations: $\hat{d} = 0.01$ (1\% duty perturbation), $\hat{\imath} = 0.04\,\si{\ampere}$ (2\% current perturbation).
  \item DC term: $D I = 0.5 \times 2.0 = 1.0\,\si{\ampere}$.
  \item First-order AC terms: $D \hat{\imath} = 0.5 \times 0.04 = 0.02\,\si{\ampere}$, and $I \hat{d} = 2.0 \times 0.01 = 0.02\,\si{\ampere}$.
  \item Total first-order perturbation: $D \hat{\imath} + I \hat{d} = 0.04\,\si{\ampere}$.
  \item Second-order product term: $\hat{d} \hat{\imath} = 0.01 \times 0.04 = 0.0004\,\si{\ampere}$.
\end{itemize}
The second-order product $\hat{d}\hat{\imath}$ is exactly $1\%$ of the first-order variation ($0.0004 / 0.04 = 0.01$). Hence, neglecting second-order products ($\hat{d}\hat{x} \approx 0$, $\hat{d}\hat{v}_{in} \approx 0$) introduces negligible error while converting a non-linear system into a solvable linear system.

\paragraph{Step 6: Linearization \& Separation into DC and AC Equations}
Substitute the perturbed variables into the averaged state equation and expand:
\begin{equation}
\begin{aligned}
\frac{d(X + \hat{x})}{dt} = &\; \left[ (D + \hat{d}) A_1 + (1 - D - \hat{d}) A_2 \right] (X + \hat{x}) \\
&+ \left[ (D + \hat{d}) B_1 + (1 - D - \hat{d}) B_2 \right] (V_{in} + \hat{v}_{in})
\end{aligned}
\end{equation}
Expanding and collecting terms of identical order:
\begin{enumerate}
  \item \textbf{DC Steady-State Equation (Zero-Order Terms):}
  Since $X$ is constant in steady state ($\dot{X} = 0$):
  \begin{equation}
  0 = \left[ D A_1 + (1-D) A_2 \right] X + \left[ D B_1 + (1-D) B_2 \right] V_{in} = A X + B V_{in}
  \end{equation}
  Solving for the DC quiescent operating point:
  \begin{equation}
  \boxed{X = -A^{-1} B V_{in}}
  \end{equation}
  where:
  \begin{equation}
  A \triangleq D A_1 + (1-D) A_2, \qquad B \triangleq D B_1 + (1-D) B_2
  \end{equation}

  \item \textbf{Small-Signal AC Dynamic Equation (First-Order Terms):}
  Subtracting the DC equilibrium equation and discarding second-order products ($\hat{d}\hat{x} \approx 0$, $\hat{d}\hat{v}_{in} \approx 0$):
  \begin{equation}
  \boxed{\dot{\hat{x}}(t) = A \hat{x}(t) + F \hat{d}(t) + B \hat{v}_{in}(t)}
  \end{equation}
  where the duty-ratio perturbation matrix $F \in \mathbb{R}^{2\times 1}$ is:
  \begin{equation}
  \boxed{F \triangleq (A_1 - A_2) X + (B_1 - B_2) V_{in}}
  \end{equation}
  The dynamic system is now completely linear and time-invariant (LTI)!
\end{enumerate}

---

\subsection{Laplace Transfer Function Matrix \& Small-Signal Network}

\subsubsection{Derivation of the Resolvent Matrix Equation}
Applying the unilateral Laplace transform ($\mathcal{L}\{\dot{\hat{x}}(t)\} = s \hat{x}(s)$, assuming zero initial AC state $\hat{x}(0^-) = 0$):
\begin{equation}
s \hat{x}(s) = A \hat{x}(s) + F \hat{d}(s) + B \hat{v}_{in}(s)
\end{equation}
Rearranging to group state variables on the left-hand side:
\begin{equation}
(s I - A) \hat{x}(s) = F \hat{d}(s) + B \hat{v}_{in}(s)
\end{equation}
where $I = \begin{bmatrix} 1 & 0 \\ 0 & 1 \end{bmatrix}$ is the $2\times 2$ identity matrix. Multiplying both sides by the inverse matrix $(sI - A)^{-1}$ (the \textbf{resolvent matrix}):
\begin{equation}
\boxed{\hat{x}(s) = (s I - A)^{-1} F \hat{d}(s) + (s I - A)^{-1} B \hat{v}_{in}(s)}
\end{equation}

\subsubsection{Resolvent Matrix Formulation for $2\times 2$ Systems}
The inverse of any $2\times 2$ matrix $M = \begin{bmatrix} a & b \\ c & e \end{bmatrix}$ is given by:
\begin{equation}
M^{-1} = \frac{\text{adj}(M)}{\det(M)} = \frac{1}{a e - b c} \begin{bmatrix} e & -b \\ -c & a \end{bmatrix}
\end{equation}
Applying this to the characteristic matrix $(s I - A)$:
\begin{equation}
(s I - A)^{-1} = \frac{\text{adj}(s I - A)}{\det(s I - A)}
\end{equation}
\begin{itemize}
  \item \textbf{Characteristic Equation:} $\det(s I - A) = 0$. The roots of $\det(s I - A) = 0$ define the open-loop poles of the switching converter.
  \item For all passive second-order $LC$ converters without internal right-half-plane poles, $\det(s I - A)$ forms a stable second-order polynomial:
  \begin{equation}
  \det(s I - A) = s^2 + 2\zeta \omega_0 s + \omega_0^2 = s^2 + \frac{\omega_0}{Q} s + \omega_0^2
  \end{equation}
\end{itemize}

\subsubsection{Standard Transfer Function Definitions}
From the state vector $\hat{x}(s) = \begin{bmatrix} \hat{\imath}_L(s) \\ \hat{v}_C(s) \end{bmatrix}$ and output $\hat{v}_o(s) = \hat{v}_C(s)$, we define:
\begin{enumerate}
  \item \textbf{Control-to-Output Transfer Function $G_{vd}(s)$:}
  Evaluated with line voltage held constant ($\hat{v}_{in}(s) = 0$):
  \begin{equation}
  \boxed{G_{vd}(s) \triangleq \left.\frac{\hat{v}_o(s)}{\hat{d}(s)}\right|_{\hat{v}_{in}=0} = \begin{bmatrix} 0 & 1 \end{bmatrix} (s I - A)^{-1} F}
  \end{equation}
  This is the primary plant transfer function required to design the voltage regulation feedback loop.

  \item \textbf{Line-to-Output Transfer Function (Audio Susceptibility) $G_{vg}(s)$:}
  Evaluated with control duty cycle held constant ($\hat{d}(s) = 0$):
  \begin{equation}
  \boxed{G_{vg}(s) \triangleq \left.\frac{\hat{v}_o(s)}{\hat{v}_{in}(s)}\right|_{\hat{d}=0} = \begin{bmatrix} 0 & 1 \end{bmatrix} (s I - A)^{-1} B}
  \end{equation}
  Quantifies how input supply ripple (e.g., $100\,\si{\hertz}$ rectified ripple) propagates to the DC output.

  \item \textbf{Control-to-Inductor Current Transfer Function $G_{id}(s)$:}
  Evaluated with $\hat{v}_{in}(s) = 0$:
  \begin{equation}
  \boxed{G_{id}(s) \triangleq \left.\frac{\hat{\imath}_L(s)}{\hat{d}(s)}\right|_{\hat{v}_{in}=0} = \begin{bmatrix} 1 & 0 \end{bmatrix} (s I - A)^{-1} F}
  \end{equation}
  Essential for designing the inner current loop in Peak Current-Mode (PCM) or Average Current-Mode (ACM) control.

  \item \textbf{Open-Loop Output Impedance $Z_{out}(s)$:}
  With perturbations applied to load current $\hat{\imath}_o(s)$ while $\hat{d}=0, \hat{v}_{in}=0$:
  \begin{equation}
  \boxed{Z_{out}(s) \triangleq \left.-\frac{\hat{v}_o(s)}{\hat{\imath}_o(s)}\right|_{\hat{d}=0, \hat{v}_{in}=0}}
  \end{equation}
  Governs transient output voltage drop during abrupt load steps ($\Delta I_o$).
\end{enumerate}

---

\subsection{Buck Converter: Complete State-Space Modeling \& Transfer Functions}

\subsubsection{Circuit Schematics \& Sub-interval Equations}
Consider a standard step-down Buck converter with filter inductance $L$, output capacitance $C$, and load resistance $R$:
\begin{itemize}
  \item \textbf{State variables:} $x(t) = \begin{bmatrix} i_L(t) \\ v_C(t) \end{bmatrix}$. Output voltage $v_o(t) = v_C(t)$.
  \item \textbf{Sub-interval 1 (Switch ON, $0 < t < d T_s$):}
  Transistor $Q$ is ON, diode $D$ is reverse biased.
  \begin{equation}
  L \frac{di_L}{dt} = v_{in} - v_C \implies \frac{di_L}{dt} = -\frac{1}{L} v_C + \frac{1}{L} v_{in}
  \end{equation}
  \begin{equation}
  C \frac{dv_C}{dt} = i_L - \frac{v_C}{R} \implies \frac{dv_C}{dt} = \frac{1}{C} i_L - \frac{1}{RC} v_C
  \end{equation}
  In matrix form:
  \begin{equation}
  A_1 = \begin{bmatrix} 0 & -1/L \\ 1/C & -1/(RC) \end{bmatrix}, \qquad B_1 = \begin{bmatrix} 1/L \\ 0 \end{bmatrix}
  \end{equation}

  \item \textbf{Sub-interval 2 (Switch OFF, $d T_s < t < T_s$):}
  Transistor $Q$ is OFF, diode $D$ freewheels.
  \begin{equation}
  L \frac{di_L}{dt} = -v_C \implies \frac{di_L}{dt} = -\frac{1}{L} v_C
  \end{equation}
  \begin{equation}
  C \frac{dv_C}{dt} = i_L - \frac{v_C}{R} \implies \frac{dv_C}{dt} = \frac{1}{C} i_L - \frac{1}{RC} v_C
  \end{equation}
  In matrix form:
  \begin{equation}
  A_2 = \begin{bmatrix} 0 & -1/L \\ 1/C & -1/(RC) \end{bmatrix}, \qquad B_2 = \begin{bmatrix} 0 \\ 0 \end{bmatrix}
  \end{equation}
\end{itemize}

\subsubsection{Averaged Matrices \& DC Steady-State}
Notice that for the Buck converter, $A_1 = A_2$. Therefore:
\begin{equation}
A = D A_1 + (1-D) A_2 = A_1 = \begin{bmatrix} 0 & -1/L \\ 1/C & -1/(RC) \end{bmatrix}
\end{equation}
\begin{equation}
B = D B_1 + (1-D) B_2 = \begin{bmatrix} D/L \\ 0 \end{bmatrix}
\end{equation}
The perturbation vector $F$ is:
\begin{equation}
F = (A_1 - A_2) X + (B_1 - B_2) V_{in} = \begin{bmatrix} 0 & 0 \\ 0 & 0 \end{bmatrix} X + \begin{bmatrix} 1/L \\ 0 \end{bmatrix} V_{in} = \begin{bmatrix} V_{in}/L \\ 0 \end{bmatrix}
\end{equation}

\paragraph{DC Operating Point Calculation:}
\begin{equation}
\det(A) = 0 - \left(-\frac{1}{LC}\right) = \frac{1}{LC}
\end{equation}
\begin{equation}
A^{-1} = L C \begin{bmatrix} -\frac{1}{RC} & \frac{1}{L} \\ -\frac{1}{C} & 0 \end{bmatrix} = \begin{bmatrix} -\frac{L}{R} & C \\ -L & 0 \end{bmatrix}
\end{equation}
\begin{equation}
X = -A^{-1} B V_{in} = -\begin{bmatrix} -\frac{L}{R} & C \\ -L & 0 \end{bmatrix} \begin{bmatrix} \frac{D V_{in}}{L} \\ 0 \end{bmatrix} = -\begin{bmatrix} -\frac{D V_{in}}{R} \\ -D V_{in} \end{bmatrix} = \begin{bmatrix} \frac{D V_{in}}{R} \\ D V_{in} \end{bmatrix} = \begin{bmatrix} I_L \\ V_o \end{bmatrix}
\end{equation}
This perfectly reproduces the well-known DC relations: $V_o = D V_{in}$ and $I_L = I_o = V_o / R$.

\subsubsection{Resolvent Matrix \& Transfer Function Derivations}
Form the characteristic matrix $(s I - A)$:
\begin{equation}
(s I - A) = \begin{bmatrix} s & 1/L \\ -1/C & s + \frac{1}{RC} \end{bmatrix}
\end{equation}
The determinant is:
\begin{equation}
\det(s I - A) = s\left(s + \frac{1}{RC}\right) + \frac{1}{LC} = s^2 + \frac{1}{RC} s + \frac{1}{LC}
\end{equation}
The adjugate matrix is:
\begin{equation}
\text{adj}(s I - A) = \begin{bmatrix} s + \frac{1}{RC} & -1/L \\ 1/C & s \end{bmatrix}
\end{equation}

\paragraph{1. Control-to-Output Transfer Function $G_{vd}(s)$:}
\begin{equation}
\begin{aligned}
\hat{x}(s) &= \frac{1}{s^2 + \frac{1}{RC} s + \frac{1}{LC}} \begin{bmatrix} s + \frac{1}{RC} & -1/L \\ 1/C & s \end{bmatrix} \begin{bmatrix} V_{in}/L \\ 0 \end{bmatrix} \hat{d}(s) \\
&= \frac{1}{s^2 + \frac{1}{RC} s + \frac{1}{LC}} \begin{bmatrix} \left(s + \frac{1}{RC}\right)\frac{V_{in}}{L} \\ \frac{V_{in}}{LC} \end{bmatrix} \hat{d}(s)
\end{aligned}
\end{equation}
Extracting the second component ($\hat{v}_o(s) = \hat{v}_C(s)$):
\begin{equation}
G_{vd}(s) = \frac{\hat{v}_o(s)}{\hat{d}(s)} = \frac{\frac{V_{in}}{LC}}{s^2 + \frac{1}{RC} s + \frac{1}{LC}} = \frac{V_{in}}{L C s^2 + \frac{L}{R} s + 1}
\end{equation}
Writing in standard canonical second-order form:
\begin{equation}
\boxed{G_{vd}(s) = \frac{G_{vd0}}{1 + \dfrac{s}{Q \omega_0} + \dfrac{s^2}{\omega_0^2}} = \frac{V_{in}}{1 + \dfrac{2\zeta}{\omega_0} s + \dfrac{s^2}{\omega_0^2}}}
\end{equation}
where:
\begin{equation}
\boxed{\omega_0 = \frac{1}{\sqrt{LC}}, \qquad Q = R\sqrt{\frac{C}{L}}, \qquad \zeta = \frac{1}{2Q} = \frac{1}{2R}\sqrt{\frac{L}{C}}, \qquad G_{vd0} = V_{in}}
\end{equation}

\paragraph{2. Control-to-Inductor Current Transfer Function $G_{id}(s)$:}
Extracting the first state component ($\hat{\imath}_L(s)$):
\begin{equation}
\boxed{G_{id}(s) = \frac{\hat{\imath}_L(s)}{\hat{d}(s)} = \frac{\frac{V_{in}}{L}\left(s + \frac{1}{RC}\right)}{s^2 + \frac{1}{RC}s + \frac{1}{LC}} = \frac{V_{in}}{R} \frac{1 + sRC}{LC s^2 + \frac{L}{R} s + 1}}
\end{equation}
Notice $G_{id}(s)$ has a Left-Half-Plane (LHP) zero at $\omega_z = -1/(RC)$ which provides phase lead ($+90^\circ$).

\paragraph{3. Line-to-Output (Audio Susceptibility) Transfer Function $G_{vg}(s)$:}
\begin{equation}
\boxed{G_{vg}(s) = \frac{\hat{v}_o(s)}{\hat{v}_{in}(s)} = \frac{\frac{D}{LC}}{s^2 + \frac{1}{RC}s + \frac{1}{LC}} = \frac{D}{LC s^2 + \frac{L}{R} s + 1}}
\end{equation}
At DC ($s = 0$), $G_{vg}(0) = D = V_o / V_{in}$, as required.

\paragraph{Physical Insight (Minimum-Phase System):}
The Buck converter control-to-output transfer function $G_{vd}(s)$ has no finite zeros (it is an all-pole second-order system). It exhibits minimum-phase behavior: the phase lag starts at $0^\circ$ at low frequency and approaches $-180^\circ$ at high frequency. It can easily be stabilized with a standard Type-II or Type-III compensator, and the loop crossover frequency $f_c$ can be designed up to $f_s / 5$.
\begin{itemize}
  \item \textbf{Capacitor Equivalent Series Resistance (ESR):} In practical capacitors with series resistance $r_C$, an LHP zero is introduced:
  \begin{equation}
  \omega_{z,ESR} = \frac{1}{r_C C}
  \end{equation}
  Because this zero lies in the left half-plane ($\mathrm{Re}\{s\} < 0$), it boosts phase lead by $+90^\circ$ and aids stability, though it reduces high-frequency ripple attenuation.
\end{itemize}

---

\subsection{Boost Converter: Complete State-Space Modeling \& Transfer Functions}

\subsubsection{Circuit Schematics \& Sub-interval Equations}
Consider a Boost step-up converter with inductor $L$, switch $Q$, diode $D$, capacitor $C$, and load resistor $R$:
\begin{itemize}
  \item \textbf{Sub-interval 1 (Switch ON, $0 < t < d T_s$):}
  Transistor $Q$ conducts. Diode $D$ is reverse biased by $v_C$.
  \begin{equation}
  L \frac{di_L}{dt} = v_{in} \implies \frac{di_L}{dt} = \frac{1}{L} v_{in}
  \end{equation}
  \begin{equation}
  C \frac{dv_C}{dt} = -\frac{v_C}{R} \implies \frac{dv_C}{dt} = -\frac{1}{RC} v_C
  \end{equation}
  In matrix form:
  \begin{equation}
  A_1 = \begin{bmatrix} 0 & 0 \\ 0 & -1/(RC) \end{bmatrix}, \qquad B_1 = \begin{bmatrix} 1/L \\ 0 \end{bmatrix}
  \end{equation}

  \item \textbf{Sub-interval 2 (Switch OFF, $d T_s < t < T_s$):}
  Transistor $Q$ is OFF. Inductor current freewheels through diode $D$ to capacitor $C$ and load $R$.
  \begin{equation}
  L \frac{di_L}{dt} = v_{in} - v_C \implies \frac{di_L}{dt} = -\frac{1}{L} v_C + \frac{1}{L} v_{in}
  \end{equation}
  \begin{equation}
  C \frac{dv_C}{dt} = i_L - \frac{v_C}{R} \implies \frac{dv_C}{dt} = \frac{1}{C} i_L - \frac{1}{RC} v_C
  \end{equation}
  In matrix form:
  \begin{equation}
  A_2 = \begin{bmatrix} 0 & -1/L \\ 1/C & -1/(RC) \end{bmatrix}, \qquad B_2 = \begin{bmatrix} 1/L \\ 0 \end{bmatrix}
  \end{equation}
\end{itemize}

\subsubsection{Averaged Matrices \& DC Steady-State Derivation}
Averaging over the switching period with $D' \triangleq 1 - D$:
\begin{equation}
A = D A_1 + (1-D) A_2 = \begin{bmatrix} 0 & -\frac{1-D}{L} \\ \frac{1-D}{C} & -\frac{1}{RC} \end{bmatrix} = \begin{bmatrix} 0 & -D'/L \\ D'/C & -1/(RC) \end{bmatrix}
\end{equation}
\begin{equation}
B = D B_1 + (1-D) B_2 = \begin{bmatrix} 1/L \\ 0 \end{bmatrix}
\end{equation}
The difference matrices are:
\begin{equation}
A_1 - A_2 = \begin{bmatrix} 0 & 1/L \\ -1/C & 0 \end{bmatrix}, \qquad B_1 - B_2 = \begin{bmatrix} 0 \\ 0 \end{bmatrix}
\end{equation}
The dynamic perturbation vector $F$ is:
\begin{equation}
\boxed{F = (A_1 - A_2) X + (B_1 - B_2) V_{in} = \begin{bmatrix} 0 & 1/L \\ -1/C & 0 \end{bmatrix} \begin{bmatrix} I_L \\ V_o \end{bmatrix} = \begin{bmatrix} V_o/L \\ -I_L/C \end{bmatrix}}
\end{equation}

\paragraph{DC Steady-State Quiescent Operating Point $X$:}
\begin{equation}
\det(A) = 0 - \left( -\frac{1-D}{L} \cdot \frac{1-D}{C} \right) = \frac{(1-D)^2}{LC} = \frac{{D'}^2}{LC}
\end{equation}
The inverse matrix $A^{-1}$ is:
\begin{equation}
A^{-1} = \frac{LC}{(1-D)^2} \begin{bmatrix} -\frac{1}{RC} & \frac{1-D}{L} \\ -\frac{1-D}{C} & 0 \end{bmatrix}
\end{equation}
Computing $X = -A^{-1} B V_{in}$:
\begin{equation}
\begin{aligned}
X &= -\frac{LC}{(1-D)^2} \begin{bmatrix} -\frac{1}{RC} & \frac{1-D}{L} \\ -\frac{1-D}{C} & 0 \end{bmatrix} \begin{bmatrix} \frac{V_{in}}{L} \\ 0 \end{bmatrix} \\
&= -\frac{LC}{(1-D)^2} \begin{bmatrix} -\frac{V_{in}}{RCL} \\ -\frac{(1-D)V_{in}}{LC} \end{bmatrix} = \begin{bmatrix} \frac{V_{in}}{(1-D)^2 R} \\ \frac{V_{in}}{1-D} \end{bmatrix}
\end{aligned}
\end{equation}
\begin{equation}
\boxed{V_o = \frac{V_{in}}{1-D}, \qquad I_L = \frac{V_{in}}{(1-D)^2 R} = \frac{I_o}{1-D} = \frac{V_o}{(1-D)R}}
\end{equation}
This rigorously verifies that the DC steady-state derived via matrix inversion matches the volt-second balance result.

\subsubsection{Derivation of the Boost Control-to-Output Transfer Function}
Form the characteristic matrix $(s I - A)$:
\begin{equation}
(s I - A) = \begin{bmatrix} s & \frac{1-D}{L} \\ -\frac{1-D}{C} & s + \frac{1}{RC} \end{bmatrix}
\end{equation}
The determinant is:
\begin{equation}
\det(s I - A) = s\left(s + \frac{1}{RC}\right) + \frac{(1-D)^2}{LC} = s^2 + \frac{1}{RC} s + \frac{(1-D)^2}{LC}
\end{equation}
The adjugate matrix is:
\begin{equation}
\text{adj}(s I - A) = \begin{bmatrix} s + \frac{1}{RC} & -\frac{1-D}{L} \\ \frac{1-D}{C} & s \end{bmatrix}
\end{equation}
Evaluating $\hat{x}(s) = \frac{\text{adj}(s I - A)}{\det(s I - A)} F \hat{d}(s)$:
\begin{equation}
\begin{aligned}
\begin{bmatrix} \hat{\imath}_L(s) \\ \hat{v}_o(s) \end{bmatrix} &= \frac{1}{s^2 + \frac{1}{RC} s + \frac{(1-D)^2}{LC}} \begin{bmatrix} s + \frac{1}{RC} & -\frac{1-D}{L} \\ \frac{1-D}{C} & s \end{bmatrix} \begin{bmatrix} V_o/L \\ -I_L/C \end{bmatrix} \hat{d}(s) \\
&= \frac{1}{s^2 + \frac{1}{RC} s + \frac{(1-D)^2}{LC}} \begin{bmatrix} \left(s + \frac{1}{RC}\right)\frac{V_o}{L} + \frac{1-D}{LC} I_L \\ \frac{1-D}{LC} V_o - s \frac{I_L}{C} \end{bmatrix} \hat{d}(s)
\end{aligned}
\end{equation}
Focusing on the output voltage perturbation $\hat{v}_o(s)$:
\begin{equation}
\hat{v}_o(s) = \frac{\frac{1-D}{LC} V_o - s \frac{I_L}{C}}{s^2 + \frac{1}{RC} s + \frac{(1-D)^2}{LC}} \hat{d}(s)
\end{equation}
Substitute the DC relation $I_L = \frac{V_o}{(1-D)R}$ into the numerator:
\begin{equation}
\text{Numerator} = \frac{1-D}{LC} V_o - s \frac{V_o}{(1-D)RC} = \frac{1-D}{LC} V_o \left[ 1 - s \frac{L}{(1-D)^2 R} \right]
\end{equation}
Factor out $\frac{(1-D)^2}{LC}$ from the denominator:
\begin{equation}
\text{Denominator} = \frac{(1-D)^2}{LC} \left[ 1 + \frac{L}{(1-D)^2 R} s + \frac{LC}{(1-D)^2} s^2 \right]
\end{equation}
Dividing numerator by denominator yields the exact control-to-output transfer function:
\begin{equation}
\boxed{G_{vd}(s) = \frac{\hat{v}_o(s)}{\hat{d}(s)} = \frac{V_o}{1-D} \frac{1 - s \dfrac{L}{(1-D)^2 R}}{1 + \dfrac{L}{(1-D)^2 R} s + \dfrac{LC}{(1-D)^2} s^2}}
\end{equation}
Or in standard second-order canonical form:
\begin{equation}
\boxed{G_{vd}(s) = G_{vd0} \frac{1 - \dfrac{s}{\omega_{z,RHP}}}{1 + \dfrac{s}{Q \omega_0} + \dfrac{s^2}{\omega_0^2}}}
\end{equation}
where:
\begin{equation}
\boxed{G_{vd0} = \frac{V_o}{1-D} = \frac{V_{in}}{(1-D)^2}}
\end{equation}
\begin{equation}
\boxed{\omega_0 = \frac{1-D}{\sqrt{LC}}, \qquad Q = (1-D) R \sqrt{\frac{C}{L}}, \qquad \zeta = \frac{1}{2Q} = \frac{1}{2(1-D)R}\sqrt{\frac{L}{C}}}
\end{equation}
\begin{equation}
\boxed{\omega_{z,RHP} = \frac{(1-D)^2 R}{L}, \qquad f_{z,RHP} = \frac{\omega_{z,RHP}}{2\pi} = \frac{(1-D)^2 R}{2\pi L}}
\end{equation}

\paragraph{Control-to-Inductor Current Transfer Function $G_{id}(s)$:}
From the first row of the matrix multiplication:
\begin{equation}
\hat{\imath}_L(s) = \frac{\left(s + \frac{1}{RC}\right)\frac{V_o}{L} + \frac{1-D}{LC} \frac{V_o}{(1-D)R}}{s^2 + \frac{1}{RC} s + \frac{(1-D)^2}{LC}} \hat{d}(s) = \frac{\frac{V_o}{L}\left(s + \frac{2}{RC}\right)}{s^2 + \frac{1}{RC} s + \frac{(1-D)^2}{LC}} \hat{d}(s)
\end{equation}
\begin{equation}
\boxed{G_{id}(s) = \frac{\hat{\imath}_L(s)}{\hat{d}(s)} = \frac{2 V_o}{(1-D)^2 R} \frac{1 + \frac{1}{2} s R C}{1 + \dfrac{L}{(1-D)^2 R} s + \dfrac{LC}{(1-D)^2} s^2}}
\end{equation}
Notice that $G_{id}(s)$ has an \textbf{LHP zero} at $s = -2/(RC)$, meaning there is \textbf{no RHP zero} in the control-to-inductor current transfer function! This is why Peak Current-Mode Control (PCMC) is so widely used for Boost converters: the inner current loop eliminates the inductor dynamics, bypassing the RHP zero within the inner loop.

---

\subsection{The Right-Half-Plane (RHP) Zero: Derivation, Physics, and Control Limits}

\subsubsection{Mathematical Origin}
Setting the numerator of $G_{vd}(s)$ to zero:
\begin{equation}
1 - s \frac{L}{(1-D)^2 R} = 0 \implies s_{z,RHP} = +\frac{(1-D)^2 R}{L} > 0
\end{equation}
Because $s_{z,RHP}$ has a positive real part ($\mathrm{Re}\{s\} > 0$), it lies strictly in the \textbf{Right Half of the complex $s$-plane} ($\sigma > 0$). Such systems are termed \textbf{Non-Minimum Phase systems}.

\subsubsection{Physical Origin: Time-Domain Transient Undershoot}
Why does a physical boost converter exhibit an RHP zero?
\begin{enumerate}
  \item Suppose the converter is in steady state, and the output voltage $v_o$ drops due to a sudden increase in load.
  \item The feedback controller responds by commanding an immediate \textbf{step increase in duty ratio} ($d \uparrow$).
  \item To increase duty ratio, the transistor switch is kept \textbf{ON for a longer duration} ($d T_s$ increases).
  \item While the transistor is ON, the diode is reverse biased ($v_D = -v_o < 0$). The inductor is connected directly across $V_{in}$ to build up magnetic energy, but it is \textbf{completely disconnected from the output stage}!
  \item During this extended ON time, the output capacitor $C$ must supply the \textbf{entire load current $I_o$ alone}:
  \begin{equation}
  \left.\frac{dv_o}{dt}\right|_{\text{switch ON}} = -\frac{V_o}{R C} < 0
  \end{equation}
  \item \textbf{The Paradoxical Undershoot:} Even though the controller commanded an increase in duty cycle to raise $v_o$, the output voltage initially \textbf{DIPS / UNDERSHOOTS} in the opposite direction!
  \item Only after the inductor current has slowly ramped up over several subsequent switching cycles and the switch finally turns OFF does the increased stored energy transfer through the diode to recharge the capacitor and eventually raise $v_o$.
\end{enumerate}
This initial negative response to a positive step input is the classic time-domain signature of an RHP zero.

\subsubsection{Frequency-Domain \& Bode Plot Behavior}
Compare the frequency response of a conventional Left-Half-Plane (LHP) zero $(1 + s/\omega_z)$ versus an RHP zero $(1 - s/\omega_z)$:
\begin{center}
\footnotesize
\renewcommand{\arraystretch}{1.2}
\begin{tabularx}{\textwidth}{@{} >{\raggedright\arraybackslash}p{3.5cm} >{\raggedright\arraybackslash}p{4.0cm} >{\raggedright\arraybackslash}p{4.0cm} >{\raggedright\arraybackslash}X @{}}
\toprule
\textbf{Zero Type} & \textbf{Laplace Factor} & \textbf{Magnitude Slope ($|G(j\omega)|$)} & \textbf{Phase Contribution ($\angle G(j\omega)$)} \\
\midrule
\textbf{Standard LHP Zero} &
$1 + \dfrac{s}{\omega_z}$ &
$+20\,\text{dB/decade}$ above $\omega_z$ &
\textbf{Phase Lead:} $+90^\circ$ \newline (improves stability margin) \\
\midrule
\textbf{Right-Half-Plane (RHP) Zero} &
$1 - \dfrac{s}{\omega_{z,RHP}}$ &
$\mathbf{+20\,\text{dB/decade}}$ above $\omega_{z,RHP}$ \newline
$\left|1 - j\frac{\omega}{\omega_z}\right| = \sqrt{1 + \left(\frac{\omega}{\omega_z}\right)^2}$ &
\textbf{Phase Lag:} $\mathbf{-90^\circ}$ \newline
$\angle\left(1 - j\frac{\omega}{\omega_z}\right) = -\arctan\left(\frac{\omega}{\omega_z}\right)$ \newline
(acts like a pole on phase!) \\
\bottomrule
\end{tabularx}
\end{center}

\begin{itemize}
  \item \textbf{Why RHP Zeros are Destructive to Control:} A normal pole provides phase lag ($-90^\circ$) but also rolls off gain ($-20\,\text{dB/dec}$), which pushes the gain below $0\,\text{dB}$ and stabilizes the loop. An RHP zero does the exact opposite: it \textbf{boosts gain by $+20\,\text{dB/dec}$} while simultaneously \textbf{injecting $-90^\circ$ of phase lag}!
  \item At frequencies above $\omega_{z,RHP}$, the combined phase lag of the two complex poles ($-180^\circ$) and the RHP zero ($-90^\circ$) reaches:
  \begin{equation}
  \text{Total Open-Loop Phase Lag} = -180^\circ - 90^\circ = -270^\circ
  \end{equation}
  This causes the Nyquist plot to encircle the critical point $(-1, j0)$, leading to catastrophic closed-loop oscillation if the loop gain crossover frequency is allowed near $\omega_{z,RHP}$.
\end{itemize}

\subsubsection{Bandwidth Limitation \& Loop Crossover Rule}
\begin{enumerate}
  \item \textbf{Forbidden Pole-Zero Cancellation:} It is impossible to cancel an RHP zero using an RHP pole in the compensator ($C(s) = \dots / (1 - s/\omega_{z,RHP})$). An unstable pole-zero cancellation in $\mathrm{Re}\{s\} > 0$ yields an internally unstable feedback system that diverges under any infinitesimal parameter mismatch or disturbance.
  \item \textbf{Fundamental Crossover Frequency Constraint:}
  To preserve adequate closed-loop phase margin ($\mathrm{PM} \ge 45^\circ - 60^\circ$), the loop gain crossover frequency $f_c$ (gain bandwidth) must be constrained well below the RHP zero frequency:
  \begin{equation}
  \boxed{f_c < \frac{1}{3 \sim 5} f_{z,RHP} = \frac{(1-D)^2 R}{2\pi (3 \sim 5) L}}
  \end{equation}
  \item \textbf{Worst-Case Operating Point Analysis:}
  The RHP zero frequency is a strong function of the operating point ($D$ and $R$):
  \begin{equation}
  f_{z,RHP} = \frac{(1-D)^2 R}{2\pi L}
  \end{equation}
  The RHP zero moves to its lowest (most restrictive) frequency under:
  \begin{itemize}
    \item \textbf{Maximum Duty Cycle $D_{\max}$:} Occurs at minimum input voltage $V_{in,\min}$ ($1 - D = V_{in}/V_o$).
    \item \textbf{Minimum Load Resistance $R_{\min}$:} Occurs at maximum load current / maximum output power ($P_{o,\max} = V_o^2 / R_{\min}$).
  \end{itemize}
  Therefore, the feedback loop must be compensated for the \textbf{worst-case minimum RHP zero}:
  \begin{equation}
  \boxed{f_{z,RHP,\min} = \frac{(1 - D_{\max})^2 R_{\min}}{2\pi L} = \left(\frac{V_{in,\min}}{V_o}\right)^2 \frac{R_{\min}}{2\pi L}}
  \end{equation}
  \item \textbf{Design Trade-Off:} Decreasing the inductance $L$ pushes $f_{z,RHP}$ to higher frequencies, allowing faster control bandwidth, but increases the inductor peak-to-peak current ripple $\Delta i_L$. Increasing $C$ helps hold up $v_o$ during the switch ON phase but does not eliminate the RHP zero.
\end{enumerate}

---

\subsection{Buck-Boost Converter State-Space Averaged Model}

\subsubsection{Circuit Equations \& State Matrices}
In an inverting Buck-Boost converter, the energy storage inductor is switched between input $V_{in}$ and output $V_o$:
\begin{itemize}
  \item \textbf{Switch ON ($0 < t < d T_s$):} $L \frac{di_L}{dt} = v_{in}$, $C \frac{dv_C}{dt} = -\frac{v_C}{R}$.
  \begin{equation}
  A_1 = \begin{bmatrix} 0 & 0 \\ 0 & -1/(RC) \end{bmatrix}, \qquad B_1 = \begin{bmatrix} 1/L \\ 0 \end{bmatrix}
  \end{equation}
  \item \textbf{Switch OFF ($d T_s < t < T_s$):} Inductor discharges into capacitor with reversed polarity ($v_L = v_C$):
  \begin{equation}
  L \frac{di_L}{dt} = v_C, \qquad C \frac{dv_C}{dt} = -i_L - \frac{v_C}{R}
  \end{equation}
  \begin{equation}
  A_2 = \begin{bmatrix} 0 & 1/L \\ -1/C & -1/(RC) \end{bmatrix}, \qquad B_2 = \begin{bmatrix} 0 \\ 0 \end{bmatrix}
  \end{equation}
  \item \textbf{Averaged Matrices ($D' = 1 - D$):}
  \begin{equation}
  A = \begin{bmatrix} 0 & \frac{1-D}{L} \\ -\frac{1-D}{C} & -\frac{1}{RC} \end{bmatrix}, \qquad B = \begin{bmatrix} D/L \\ 0 \end{bmatrix}
  \end{equation}
  \begin{equation}
  F = (A_1 - A_2) X + (B_1 - B_2) V_{in} = \begin{bmatrix} 0 & -1/L \\ 1/C & 0 \end{bmatrix} \begin{bmatrix} I_L \\ V_o \end{bmatrix} + \begin{bmatrix} V_{in}/L \\ 0 \end{bmatrix} = \begin{bmatrix} \frac{V_{in} - V_o}{L} \\ \frac{I_L}{C} \end{bmatrix}
  \end{equation}
\end{itemize}

\subsubsection{Steady-State \& Transfer Functions}
\begin{equation}
X = -A^{-1} B V_{in} \implies V_o = -\frac{D}{1-D} V_{in}, \qquad I_L = \frac{D V_{in}}{(1-D)^2 R} = \frac{|I_o|}{1-D}
\end{equation}
The control-to-output transfer function is:
\begin{equation}
\boxed{G_{vd}(s) = \frac{V_o}{D(1-D)} \frac{1 - s \dfrac{D L}{(1-D)^2 R}}{1 + \dfrac{L}{(1-D)^2 R} s + \dfrac{LC}{(1-D)^2} s^2}}
\end{equation}
\begin{itemize}
  \item \textbf{Buck-Boost RHP Zero:}
  \begin{equation}
  \boxed{\omega_{z,RHP} = \frac{(1-D)^2 R}{D L}, \qquad f_{z,RHP} = \frac{(1-D)^2 R}{2\pi D L}}
  \end{equation}
  Notice the extra factor of $D$ in the denominator: at high duty cycles ($D > 0.5$), the Buck-Boost RHP zero drops even lower in frequency than the Boost converter zero, severely restricting transient performance.
\end{itemize}

---

\subsection{Step-by-Step Exam Recipes}

\subsubsection{Exam Recipe 1: Six-Step Procedure to Formulate SSA Matrices from Circuit Schematics}
\paragraph{Problem Description:}
Given an unfamiliar switched DC-DC topology (or Buck, Boost, SEPIC, Cuk) with switches, diodes, inductors, and capacitors, derive the state-space matrices ($A_1, B_1, A_2, B_2$), the averaged matrices ($A, B, F$), and the characteristic polynomial.

\paragraph{Step-by-Step Execution Guide:}
\begin{enumerate}
  \item \textbf{Step 1: Select State and Input Variables}
  \begin{itemize}
    \item Count the independent inductors $n_L$ and independent capacitors $n_C$. Total state dimension is $n = n_L + n_C$ (typically $n=2$).
    \item Assign state vector: $x(t) = [i_{L1}(t), \dots, v_{C1}(t), \dots]^T$.
    \item Assign input vector: $u(t) = [v_{in}(t)]^T$.
  \end{itemize}

  \item \textbf{Step 2: Draw the Switch-ON Equivalent Circuit ($0 < t < d T_s$)}
  \begin{itemize}
    \item Replace active switches with short circuits (closed) and diodes with open circuits (if reverse biased) or shorts (if forward biased).
    \item For each inductor, write KVL: $L \frac{di_L}{dt} = \sum v_k$. Express all terms strictly as linear combinations of states $x$ and inputs $u$.
    \item For each capacitor, write KCL: $C \frac{dv_C}{dt} = \sum i_k$. Express all currents as linear combinations of states $x$ and inputs $u$.
    \item Read off the coefficients to populate matrix $A_1$ (dimension $n\times n$) and $B_1$ (dimension $n\times m$).
  \end{itemize}

  \item \textbf{Step 3: Draw the Switch-OFF Equivalent Circuit ($d T_s < t < T_s$)}
  \begin{itemize}
    \item Replace active switches with open circuits and freewheeling diodes with short circuits.
    \item Repeat KVL for all inductors and KCL for all capacitors.
    \item Read off the coefficients to populate matrix $A_2$ and $B_2$.
  \end{itemize}

  \item \textbf{Step 4: Compute the Averaged Matrices}
  \begin{itemize}
    \item Average state matrix: $A = D A_1 + (1-D) A_2$.
    \item Average input matrix: $B = D B_1 + (1-D) B_2$.
    \item Check dimensions: $A$ must be $2\times 2$, $B$ must be $2\times 1$.
  \end{itemize}

  \item \textbf{Step 5: Verify DC Equilibrium $X = -A^{-1} B V_{in}$}
  \begin{itemize}
    \item Compute $\det(A)$ and adjugate $\text{adj}(A)$.
    \item Compute $X = -A^{-1} B V_{in}$.
    \item \textbf{Sanity Check:} Does $X_2 = V_o$ match the known steady-state voltage conversion ratio $M(D) V_{in}$ derived from volt-second balance? If not, check signs in Step 2 or Step 3!
  \end{itemize}

  \item \textbf{Step 6: Formulate Dynamic Perturbation Vector $F$}
  \begin{itemize}
    \item Compute difference matrices: $\Delta A = A_1 - A_2$, $\Delta B = B_1 - B_2$.
    \item Compute $F = \Delta A \cdot X + \Delta B \cdot V_{in}$.
    \item Substitute $F$ into $(s I - A)^{-1} F$ to obtain transfer functions.
  \end{itemize}
\end{enumerate}

\subsubsection{Exam Recipe 2: Calculating DC Steady-State $X = -A^{-1} B V_{in}$ and Dynamic Vector $F$}
\paragraph{Problem Description:}
Given numerical matrices $A_1, B_1, A_2, B_2$ and operating point values ($V_{in}, D, R, L, C$), compute numerical values for $X$ and $F$, and identify any sign errors.

\paragraph{Step-by-Step Procedure:}
\begin{enumerate}
  \item \textbf{Evaluate Averaged Matrix $A$:}
  \[
  A = \begin{bmatrix} a_{11} & a_{12} \\ a_{21} & a_{22} \end{bmatrix} = D \begin{bmatrix} a_{1,11} & a_{1,12} \\ a_{1,21} & a_{1,22} \end{bmatrix} + (1-D) \begin{bmatrix} a_{2,11} & a_{2,12} \\ a_{2,21} & a_{2,22} \end{bmatrix}
  \]
  \item \textbf{Calculate Determinant $\det(A)$:}
  \[
  \det(A) = a_{11} a_{22} - a_{12} a_{21}
  \]
  For passive $RLC$ converters, $\det(A) > 0$.
  \item \textbf{Compute Inverse Matrix $A^{-1}$:}
  \[
  A^{-1} = \frac{1}{\det(A)} \begin{bmatrix} a_{22} & -a_{12} \\ -a_{21} & a_{11} \end{bmatrix}
  \]
  \item \textbf{Compute $X = -A^{-1} B V_{in}$:}
  \[
  \begin{bmatrix} I_L \\ V_o \end{bmatrix} = -\frac{V_{in}}{\det(A)} \begin{bmatrix} a_{22} & -a_{12} \\ -a_{21} & a_{11} \end{bmatrix} \begin{bmatrix} b_1 \\ b_2 \end{bmatrix} = -\frac{V_{in}}{\det(A)} \begin{bmatrix} a_{22} b_1 - a_{12} b_2 \\ -a_{21} b_1 + a_{11} b_2 \end{bmatrix}
  \]
  \item \textbf{Evaluate Vector $F$:}
  \[
  F = \begin{bmatrix} f_1 \\ f_2 \end{bmatrix} = \begin{bmatrix} (a_{1,11} - a_{2,11})I_L + (a_{1,12} - a_{2,12})V_o + (b_{1,1} - b_{2,1})V_{in} \\ (a_{1,21} - a_{2,21})I_L + (a_{1,22} - a_{2,22})V_o + (b_{1,2} - b_{2,2})V_{in} \end{bmatrix}
  \]
  \item \textbf{Exam Tip on Signs:}
  \begin{itemize}
    \item In Buck: $f_1 = V_{in}/L > 0$, $f_2 = 0$.
    \item In Boost: $f_1 = V_o/L > 0$, $f_2 = -I_L/C < 0$. The \textbf{negative sign} in $f_2$ is the direct mathematical origin of the RHP zero!
  \end{itemize}
\end{enumerate}

\subsubsection{Exam Recipe 3: Boost Converter RHP Zero Sizing \& Closed-Loop Bandwidth Design}
\paragraph{Problem Statement (Comprehensive Exam Problem):}
A Boost DC-DC converter is designed for an active Power Factor Correction / renewable energy stage with the following nominal specifications:
\begin{itemize}
  \item Input voltage range: $V_{in} = 12\,\si{\volt}$ to $18\,\si{\volt}$ (nominal $V_{in} = 12\,\si{\volt}$).
  \item Regulated output voltage: $V_o = 24\,\si{\volt}$.
  \item Nominal output power: $P_o = 48\,\si{\watt}$ ($I_o = 2.0\,\si{\ampere}$, $R = 12\,\Omega$).
  \item Maximum output power: $P_{o,\max} = 96\,\si{\watt}$ ($I_{o,\max} = 4.0\,\si{\ampere}$, $R_{\min} = 6.0\,\Omega$).
  \item Filter components: $L = 100\,\mu\si{\henry}$, $C = 220\,\mu\si{\farad}$.
  \item Switching frequency: $f_s = 100\,\si{\kilo\hertz}$.
\end{itemize}

\paragraph{Questions to Solve:}
\begin{enumerate}
  \item Find the steady-state nominal duty cycle $D$ and average inductor current $I_L$.
  \item Formulate the numerical averaged state matrix $A$ and perturbation vector $F$.
  \item Calculate the natural resonance frequency $f_0$ and quality factor $Q$.
  \item Calculate the nominal RHP zero frequency $f_{z,RHP}$.
  \item Determine the worst-case minimum RHP zero frequency $f_{z,RHP,\min}$.
  \item Recommend the maximum allowable closed-loop voltage crossover frequency $f_c$.
\end{enumerate}

\paragraph{Step-by-Step Solution Procedure:}
\begin{enumerate}
  \item \textbf{Nominal DC Operating Point ($V_{in} = 12\,\si{\volt}, V_o = 24\,\si{\volt}, R = 12\,\Omega$):}
  \begin{equation}
  1 - D = \frac{V_{in}}{V_o} = \frac{12}{24} = 0.50 \implies D = 0.50
  \end{equation}
  \begin{equation}
  I_L = \frac{V_o}{(1-D)R} = \frac{24}{(1 - 0.50) \times 12} = \frac{24}{6.0} = 4.0\,\si{\ampere}
  \end{equation}
  Check: $P_{in} = V_{in} I_L = 12 \times 4.0 = 48\,\si{\watt} = P_o$.

  \item \textbf{Numerical Matrices ($A$ and $F$):}
  \begin{equation}
  \frac{1-D}{L} = \frac{0.50}{100 \times 10^{-6}} = 5000\,\si{\per\second}
  \end{equation}
  \begin{equation}
  \frac{1-D}{C} = \frac{0.50}{220 \times 10^{-6}} \approx 2272.7\,\si{\per\second}
  \end{equation}
  \begin{equation}
  \frac{1}{R C} = \frac{1}{12 \times 220 \times 10^{-6}} \approx 378.8\,\si{\per\second}
  \end{equation}
  \begin{equation}
  A = \begin{bmatrix} 0 & -5000 \\ 2272.7 & -378.8 \end{bmatrix}
  \end{equation}
  Perturbation vector $F$:
  \begin{equation}
  f_1 = \frac{V_o}{L} = \frac{24}{100 \times 10^{-6}} = 240\,000\,\si{\volt\per\henry}
  \end{equation}
  \begin{equation}
  f_2 = -\frac{I_L}{C} = -\frac{4.0}{220 \times 10^{-6}} \approx -18\,182\,\si{\ampere\per\farad}
  \end{equation}
  \begin{equation}
  F = \begin{bmatrix} 240\,000 \\ -18\,182 \end{bmatrix}
  \end{equation}

  \item \textbf{Natural Resonance Frequency $f_0$ and Quality Factor $Q$:}
  \begin{equation}
  \omega_0 = \frac{1-D}{\sqrt{LC}} = \frac{0.50}{\sqrt{100 \times 10^{-6} \times 220 \times 10^{-6}}} = \frac{0.50}{\sqrt{2.2 \times 10^{-8}}} = \frac{0.50}{1.4832 \times 10^{-4}} \approx 3371\,\si{\radian\per\second}
  \end{equation}
  \begin{equation}
  f_0 = \frac{\omega_0}{2\pi} = \frac{3371}{2\pi} \approx 536.5\,\si{\hertz}
  \end{equation}
  \begin{equation}
  Q = (1-D) R \sqrt{\frac{C}{L}} = 0.50 \times 12 \times \sqrt{\frac{220 \times 10^{-6}}{100 \times 10^{-6}}} = 6.0 \times \sqrt{2.2} \approx 6.0 \times 1.4832 \approx 8.90
  \end{equation}
  Since $Q > 0.5$ ($\zeta = 1/(2Q) \approx 0.056$), the open-loop converter is underdamped and exhibits a sharp resonant peak at $536.5\,\si{\hertz}$.

  \item \textbf{Nominal RHP Zero Frequency $f_{z,RHP}$:}
  \begin{equation}
  \omega_{z,RHP} = \frac{(1-D)^2 R}{L} = \frac{(0.50)^2 \times 12}{100 \times 10^{-6}} = \frac{0.25 \times 12}{10^{-4}} = 30\,000\,\si{\radian\per\second}
  \end{equation}
  \begin{equation}
  f_{z,RHP} = \frac{\omega_{z,RHP}}{2\pi} = \frac{30\,000}{2\pi} \approx 4774.6\,\si{\hertz} \approx 4.77\,\si{\kilo\hertz}
  \end{equation}

  \item \textbf{Worst-Case Minimum RHP Zero Frequency $f_{z,RHP,\min}$:}
  Worst case occurs at minimum input voltage $V_{in,\min} = 12\,\si{\volt}$ (maximum duty cycle $D_{\max} = 0.50$) and maximum load $P_{o,\max} = 96\,\si{\watt}$ ($R_{\min} = 6.0\,\Omega$):
  \begin{equation}
  \omega_{z,RHP,\min} = \frac{(1 - D_{\max})^2 R_{\min}}{L} = \frac{(0.50)^2 \times 6.0}{100 \times 10^{-6}} = \frac{0.25 \times 6.0}{10^{-4}} = 15\,000\,\si{\radian\per\second}
  \end{equation}
  \begin{equation}
  \boxed{f_{z,RHP,\min} = \frac{15\,000}{2\pi} \approx 2387.3\,\si{\hertz} \approx 2.39\,\si{\kilo\hertz}}
  \end{equation}
  If the input voltage drops further to $V_{in} = 8.0\,\si{\volt}$ ($1 - D = 8/24 = 1/3 \implies D = 2/3 \approx 0.667$):
  \begin{equation}
  f_{z,RHP} = \frac{(1/3)^2 \times 6.0}{2\pi \times 100 \times 10^{-6}} = \frac{(1/9) \times 6.0}{2\pi \times 10^{-4}} \approx \frac{0.6667}{6.283 \times 10^{-4}} \approx 1061\,\si{\hertz}
  \end{equation}

  \item \textbf{Recommended Loop Crossover Frequency $f_c$:} \par
  Applying the rule of thumb ($f_c \le f_{z,RHP,\min} / 5$ to $f_{z,RHP,\min} / 3$):
  \begin{equation}
  \boxed{f_c \le \frac{f_{z,RHP,\min}}{5} = \frac{2387.3\,\si{\hertz}}{5} \approx 477\,\si{\hertz}}
  \end{equation}
  \begin{itemize}
    \item \textbf{Critical Design Observation:} Notice that $f_c \approx 477\,\si{\hertz}$ is slightly below the open-loop resonance frequency $f_0 = 536.5\,\si{\hertz}$! This illustrates the severe bandwidth penalty imposed by an RHP zero in single-loop voltage-mode control.
    \item \textbf{Engineering Solution:} To break through this limit, engineers employ \textbf{Peak Current-Mode Control (PCMC)}. In PCMC, the inner current loop eliminates the inductor pole from the outer loop, moving the voltage loop crossover frequency higher while maintaining stable operation.
  \end{itemize}
\end{enumerate}
